\chapter{Hosting multiple web applications}
\label{chap:multiapps}
% Corresponds to phases 02 and 03 of the roadmap.

\section{Principle adopted}
The platform hosts each application as its own container, or its own virtual machine when the application is large or belongs to another team, and exposes it under its own subdomain of \texttt{ucac-icam.info} (for example \texttt{dashboard.ucac-icam.info} and \texttt{gitlab.ucac-icam.info}). The key property this gives is that adding a new application never requires buying a new domain: once the root domain is registered and its DNS delegated to Cloudflare, an unlimited number of subdomains can be created at no extra cost, each pointing to a different application. This mechanism, and the cost analysis behind acquiring the domain, was written up separately for the supervisors (see the standalone technical note referenced in the deployment guide) and is only summarised here.

\section{The second application: GitLab CE}
\subsection{Why a dedicated virtual machine}
GitLab is the platform on which the other interns push their code and run their tests. It is hosted on its own virtual machine, \texttt{gitlab-vm}, rather than next to the dashboard. This isolates the interns' code from Jenkins and its deployment keys, and keeps GitLab available even when \texttt{deploy-vm} is stopped for maintenance. The machine has 4 virtual CPUs, 8~GB of memory and a 50~GB disk; GitLab runs as a Docker Compose service with its data kept in folders on the host, so that the container can be recreated without losing anything. Git over SSH is exposed on port 2222, so that it never conflicts with the administration SSH on port 22, and the firewall limits the machine to the local network.

\subsection{Operations}
Two procedures were written and tested: the update (version check, free space check, backup of the data and of the secrets file, pull of the new image, wait for the end of the internal reconfiguration) and the reset of a forgotten administrator password with the \texttt{gitlab-rake} command. GitLab imposes an upgrade path between major versions, so the procedure points to the official path before any update.

\subsection{The full-disk incident}
GitLab became unreachable with an HTTP 502 error. The logs showed ``No space left on device'' errors coming from several components at once (database, Git storage, web server): a full root disk makes GitLab collapse in cascade. The dashboard's disk-usage panel, which breaks the space down by folder on the whole host (chapter~\ref{chap:supervision}), was used to find the cause without opening an SSH session. Cleaning up unused Docker data and the system journal, then restarting the GitLab services, restored the service. The incident led to a documented diagnosis-cleanup-restart procedure and to a check of the free disk space before every update.

\section{Public access and reverse proxying}
\label{sec:tunnel-details}
\subsection{One named tunnel per machine}
Each virtual machine runs its own named tunnel with its own credentials: \texttt{info-control-center} on \texttt{deploy-vm} and \texttt{info-gitlab} on \texttt{gitlab-vm}. A single tunnel for both would have been simpler, but it would have tied the availability and the secrets of one machine to the other. The configuration of a tunnel maps a public host name to an internal service, and anything else receives a 404 error. The credentials file is a secret kept out of Git, and the configuration files must exist before the containers first start, for the reason explained in chapter~\ref{chap:infra}.

\subsection{Who plays the role of reverse proxy}
Three components could play the role of reverse proxy, and it matters to know which one actually receives the public traffic. The \textbf{tunnel configuration} does: it receives the HTTPS traffic from Cloudflare and routes it to a container according to the host name. \textbf{Traefik}, deployed on \texttt{deploy-vm}, is configured by Docker labels and designed for Let's Encrypt certificates, which need ports 80 and 443 to be reachable from the Internet; behind CGNAT it is kept for a future public IP address but is not in the public path. \textbf{GitLab's own web server} listens on plain HTTP inside its container, and the GitLab tunnel reaches it directly. Adding a new application therefore means three steps: add a routing rule for its host name in the tunnel configuration, create the matching DNS route, and restart the tunnel container, with the application's container on the same Docker network.

\subsection{GitLab behind an external HTTPS}
Cloudflare terminates HTTPS, so GitLab stays in plain HTTP behind its tunnel. Its external URL is declared in \texttt{https}, its internal web server listens on port 80 without TLS, and the automatic certificate request is disabled. Two consequences were noticed. First, the web interface must be used through the domain: on the local \texttt{http} address, the session cookie is marked as secure and the browser does not send it back, so a login attempt ends with an HTTP 422 error (request rejected by the CSRF protection). Git operations over HTTP on the local network are not affected. Second, Git over SSH (port 2222) remains reachable only from the local network, because the tunnel carries web traffic only. Cloudflare's free plan also limits a request to 100~MB, so a very large \texttt{git push} must go through the local network.

\section{Network isolation between applications}
Each application runs on its own Docker network, with the tunnel container as the only component attached to more than one network -- so applications cannot reach one another directly, only through their public route. Beyond Docker networks, the strongest isolation is the one between the two virtual machines: the dashboard and Jenkins on one side, the interns' code and their pipelines on the other, each with its own firewall, its own tunnel and its own secrets.

\section{Status at the time of writing}
Two applications are deployed and reachable from the Internet under the project's domain: the supervision dashboard, which served as the platform's reference application throughout the internship, and GitLab, which is used by the other interns. Hosting a third application, and protecting these two further (see chapter~\ref{chap:bilan}), are the planned next steps.

\clearpage
